% Cover Letter: a one-page A4 application letter for academic posts and fellowships.
% Set in Libertinus Serif with one deep burgundy accent. Build with pdflatex.
% Edit the sender, recipient and subject below; the letter body follows.
\RequirePackage{fix-cm}% lets the 17 pt name scale cleanly
\documentclass[11pt]{article}
% The text block is about 120 mm wide: 65 to 70 characters per line.
\usepackage[a4paper,top=18mm,bottom=18mm,left=45mm,right=45mm]{geometry}
\usepackage[T1]{fontenc}
\usepackage[osf]{libertinus}
\usepackage[letterspace=200]{microtype}
\usepackage{xcolor}
\usepackage{tikz}
\usepackage{fontawesome5}
\usepackage[hidelinks]{hyperref}

% ---------------------------------------------------------------------------
% Sender: the letterhead and PDF metadata use these.
\newcommand{\sendername}{Tomás Ferreira Lindqvist}
\newcommand{\senderrole}{Postdoctoral Research Fellow in Palaeoclimate Modelling}
\newcommand{\senderemail}{t.lindqvist@example.org}
\newcommand{\senderphone}{+46 70 123 45 67}
\newcommand{\senderweb}{tlindqvist.example.org}
\newcommand{\senderplace}{Uppsala, Sweden}
\hypersetup{pdftitle={Letter of application -- \sendername},pdfauthor={\sendername}}

% Colours: one deep burgundy accent, a warm grey for secondary text.
\definecolor{accent}{RGB}{110,24,40}
\definecolor{quiet}{RGB}{105,98,96}
\definecolor{ink}{RGB}{30,26,26}

\setlength{\parindent}{0pt}
\setlength{\parskip}{0.62\baselineskip}
\linespread{1.08}
\pagestyle{empty}

% Contact line entry: a small accent icon before each item.
\newcommand{\contactitem}[2]{{\color{accent}\scriptsize #1}\,\,#2}

\begin{document}
\color{ink}
% ============================== LETTERHEAD ================================
% The letterhead runs wider than the body, so it frames the letter.
\noindent\makebox[\textwidth][c]{%
  \begin{minipage}[b]{\dimexpr\textwidth+50mm}
    \begin{minipage}[b]{0.64\linewidth}
      \raggedright
      {\fontsize{17}{20}\selectfont\color{accent}\textls[150]{\scshape\MakeLowercase{\sendername}}}\par
      \vspace{3pt}
      {\small\itshape\color{quiet}\senderrole}
    \end{minipage}%
    \hfill
    \begin{minipage}[b]{0.34\linewidth}
      \raggedleft\footnotesize\color{quiet}\linespread{1.25}\selectfont
      \contactitem{\faEnvelope[regular]}{\href{mailto:\senderemail}{\senderemail}}\\
      \contactitem{\faPhone}{\senderphone}\\
      \contactitem{\faGlobe}{\href{https://\senderweb}{\senderweb}}\\
      \contactitem{\faMapMarker*}{\senderplace}
    \end{minipage}\par
    \vspace{6pt}
    {\color{accent}\rule{\linewidth}{0.5pt}}
  \end{minipage}}
\par\vspace{10mm}

% ============================== RECIPIENT AND DATE =========================
\begin{minipage}[t]{0.62\textwidth}
  \linespread{1.12}\selectfont
  Professor Ingrid Halvorsen\\
  Chair, Search Committee\\
  Department of Earth Sciences\\
  University of Bergen\\
  Allégaten 41, 5007 Bergen, Norway
\end{minipage}%
\hfill
\begin{minipage}[t]{0.3\textwidth}
  \raggedleft 3 March 2026
\end{minipage}
\par\vspace{7mm}

% Subject line: letterspaced small caps in the accent colour.
{\color{accent}\textls[60]{\scshape re: associate professor in climate dynamics, ref.\ 2026/1184}}
\par\vspace{2mm}

% ============================== BODY =====================================
Dear Professor Halvorsen,

I am writing to apply for the Associate Professorship in Climate
Dynamics advertised by the Department of Earth Sciences. I am a
postdoctoral fellow at the Bolin Centre in Uppsala, where I study how
the Atlantic overturning circulation responded to abrupt warming during
the last deglaciation, and how those events can constrain the models we
use to project the coming century.

My research joins three strands that rarely meet in one person: proxy
records from marine sediment cores, idealised ocean models, and the
statistical tools that link the two. In the past four years this work
has produced eleven peer-reviewed papers, including a 2025 article in
\emph{Nature Geoscience} showing that two rival reconstructions of the
Younger Dryas are consistent once seasonal bias in the proxies is
modelled explicitly. Bergen's strengths in ocean observation and the
Bjerknes Centre's modelling groups are exactly the setting in which
this programme can grow.

Teaching is a part of the job I look forward to. I have designed and
taught a graduate course on palaeoclimate data analysis, supervised two
master's theses to completion, and co-supervise a doctoral student. I
would be glad to contribute to the department's courses in physical
oceanography and to build a seminar that brings model and proxy
students into the same room.

I enclose my curriculum vitae, a research statement, a teaching
statement and three representative publications. Thank you for your
time and consideration; I would welcome the chance to discuss how my
work could support the department's plans.

% ============================== CLOSING ==================================
\vspace{1mm}
Yours sincerely,\par
\vspace{1mm}
% Signature placeholder: a TikZ squiggle. Replace with a scanned signature
% via \includegraphics[height=12mm]{signature} if you have one.
\begin{tikzpicture}[x=1mm,y=1mm,line cap=round,line join=round]
  \draw[ink!85,line width=0.75pt] plot[smooth,tension=0.7] coordinates
    {(0,2) (2,8) (4,3) (5,0) (6,5) (9,9) (10,4) (11,1) (14,6) (16,3)
     (19,5) (22,2) (25,6) (27,3) (31,5) (36,4)};
  \draw[ink!85,line width=0.6pt] plot[smooth,tension=0.8] coordinates
    {(3,-1.5) (15,-0.5) (28,-1) (40,0.5)};
\end{tikzpicture}\par
\vspace{1mm}
\sendername

\vfill
{\footnotesize\color{quiet}\textls[60]{\scshape enclosures}\enspace
Curriculum vitae; research statement; teaching statement; three publications.}
\end{document}
